\section{Przykłady działania programu}

\noindent \textbf{Przykład 1: Multigraf bez pętli i z cyklem Hamiltona} \\

\noindent \textbf{Argumenty wejściowe:}
\begin{itemize}
    \item Liczba grafów: \(1\)
    \item Liczba wierzchołków w grafie: \(3\)
    \item Macierz sąsiedztwa:
    \[
    A = \begin{bmatrix}
    0 & 2 & 1 \\
    2 & 0 & 1 \\
    1 & 1 & 0 \\
    \end{bmatrix}
    \]
    gdzie \(A[i][j]\) oznacza liczbę krawędzi pomiędzy wierzchołkiem \(i\) a \(j\).
\end{itemize}

\noindent \textbf{Wynik:}
\begin{itemize}
    \item Rozmiar grafu (liczba krawędzi): \(4\)
    \item Wszystkie cykle: \(0 \to 1 \to 2 \to 0\)
    \item Cykl Hamiltona: Tak (\(0 \to 1 \to 2 \to 0\))
    \item Minimalne rozszerzenie dla cyklu Hamiltona: \(0\)
    \item Maksymalna długość cyklu: \(3\)
\end{itemize}

\vspace{1em}

\noindent \textbf{Przykład 2: Multigraf z pętlami i cyklami różnej długości} \\

\noindent \textbf{Argumenty wejściowe:}
\begin{itemize}
    \item Liczba grafów: \(2\)
    \item Graf 1: \(3\) wierzchołki, macierz sąsiedztwa:
    \[
    A_1 = \begin{bmatrix}
    0 & 1 & 1 \\
    1 & 0 & 0 \\
    1 & 0 & 0 \\
    \end{bmatrix}
    \]
    \item Graf 2: \(4\) wierzchołki, macierz sąsiedztwa:
    \[
    A_2 = \begin{bmatrix}
    0 & 2 & 0 & 1 \\
    2 & 0 & 1 & 0 \\
    0 & 1 & 0 & 1 \\
    1 & 0 & 1 & 0 \\
    \end{bmatrix}
    \]
\end{itemize}

\noindent \textbf{Wynik:}
\begin{itemize}
    \item \textbf{Graf 1:}
    \begin{itemize}
        \item Rozmiar grafu: \(2\)
        \item Wszystkie cykle: Brak
        \item Cykl Hamiltona: Nie
        \item Minimalne rozszerzenie dla cyklu Hamiltona: \(1\)
        \item Maksymalna długość cyklu: \(2\)
    \end{itemize}
    \item \textbf{Graf 2:}
    \begin{itemize}
        \item Rozmiar grafu: \(5\)
        \item Wszystkie cykle: \(0 \to 1 \to 2 \to 3 \to 0\)
        \item Cykl Hamiltona: Tak (\(0 \to 1 \to 2 \to 3 \to 0\))
        \item Minimalne rozszerzenie dla cyklu Hamiltona: \(0\)
        \item Maksymalna długość cyklu: \(4\)
    \end{itemize}
    \item Metryka podobieństwa pomiędzy grafem 1 i grafem 2: \(1.167\)
\end{itemize}

\vspace{1em}

\noindent \textbf{Przykład 3: Minimalne rozszerzenie dla cyklu Hamiltona} \\

\noindent \textbf{Argumenty wejściowe:}
\begin{itemize}
    \item Liczba grafów: \(3\)
    \item Graf 1: \(3\) wierzchołki, macierz sąsiedztwa:
    \[
    A_1 = \begin{bmatrix}
    0 & 2 & 1 \\
    2 & 0 & 1 \\
    1 & 1 & 0 \\
    \end{bmatrix}
    \]
    \item Graf 2: \(4\) wierzchołki, macierz sąsiedztwa:
    \[
    A_2 = \begin{bmatrix}
    0 & 3 & 1 & 0 \\
    3 & 0 & 0 & 2 \\
    1 & 0 & 0 & 1 \\
    0 & 2 & 1 & 0 \\
    \end{bmatrix}
    \]
    \item Graf 3: \(4\) wierzchołki, macierz sąsiedztwa:
    \[
    A_3 = \begin{bmatrix}
    0 & 2 & 0 & 0 \\
    2 & 0 & 1 & 2 \\
    0 & 1 & 0 & 1 \\
    0 & 2 & 1 & 0 \\
    \end{bmatrix}
    \]
\end{itemize}

\noindent \textbf{Wynik:}
\begin{itemize}
    \item \textbf{Graf 1:}
    \begin{itemize}
        \item Rozmiar grafu: \(4\)
        \item Wszystkie cykle: \(0 \to 1 \to 2 \to 0\)
        \item Cykl Hamiltona: Tak (\(0 \to 1 \to 2 \to 0\))
        \item Minimalne rozszerzenie dla cyklu Hamiltona: \(0\)
        \item Maksymalna długość cyklu: \(3\)
    \end{itemize}
    \item \textbf{Graf 2:}
    \begin{itemize}
        \item Rozmiar grafu: \(7\)
        \item Wszystkie cykle: \(0 \to 1 \to 3 \to 2 \to 0\)
        \item Cykl Hamiltona: Tak (\(0 \to 1 \to 3 \to 2 \to 0\))
        \item Minimalne rozszerzenie dla cyklu Hamiltona: \(0\)
        \item Maksymalna długość cyklu: \(4\)
    \end{itemize}
    \item \textbf{Graf 3:}
    \begin{itemize}
        \item Rozmiar grafu: \(6\)
        \item Wszystkie cykle: \(1 \to 2 \to 3 \to 1\)
        \item Cykl Hamiltona: Nie
        \item Minimalne rozszerzenie dla cyklu Hamiltona: \(1\)
        \item Maksymalna długość cyklu: \(3\)
    \end{itemize}
    \item Metryka podobieństwa pomiędzy grafem 1 i grafem 2: \(0.833\)
    \item Metryka podobieństwa pomiędzy grafem 2 i grafem 3: \(0.500\)
\end{itemize}

\vspace{1em}

\noindent \textbf{Przykład 4: Metryka w przestrzeni multigrafów} \\

\noindent \textbf{Argumenty wejściowe:}
\begin{itemize}
    \item Liczba grafów: \(3\)
    \item Graf 1: \(3\) wierzchołki, macierz sąsiedztwa:
    \[
    A_1 = \begin{bmatrix}
    0 & 0 & 0 \\
    0 & 0 & 0 \\
    0 & 0 & 0 \\
    \end{bmatrix}
    \]
    \item Graf 2: \(2\) wierzchołki, macierz sąsiedztwa:
    \[
    A_2 = \begin{bmatrix}
    0 & 1 \\
    1 & 0 \\
    \end{bmatrix}
    \]
    \item Graf 3: \(5\) wierzchołków, macierz sąsiedztwa:
    \[
    A_3 = \begin{bmatrix}
    0 & 1 & 1 & 1 & 0 \\
    1 & 0 & 0 & 2 & 1 \\
    1 & 0 & 0 & 0 & 1 \\
    1 & 2 & 0 & 0 & 1 \\
    0 & 1 & 1 & 1 & 0 \\
    \end{bmatrix}
    \]
\end{itemize}

\noindent \textbf{Wynik:}
\begin{itemize}
    \item \textbf{Graf 1:}
    \begin{itemize}
        \item Rozmiar grafu: \(0\)
        \item Wszystkie cykle: Brak
        \item Cykl Hamiltona: Nie
        \item Minimalne rozszerzenie dla cyklu Hamiltona: \(3\)
        \item Maksymalna długość cyklu: \(0\)
    \end{itemize}
    \item \textbf{Graf 2:}
    \begin{itemize}
        \item Rozmiar grafu: \(1\)
        \item Wszystkie cykle: Brak
        \item Cykl Hamiltona: Tak (\(0 \to 1 \to 0\))
        \item Minimalne rozszerzenie dla cyklu Hamiltona: \(0\)
        \item Maksymalna długość cyklu: \(2\)
    \end{itemize}
    \item \textbf{Graf 3:}
    \begin{itemize}
        \item Rozmiar grafu: \(8\)
        \item Wszystkie cykle: \(6\)
        \item Cykl Hamiltona: Tak (\(0 \to 1 \to 3 \to 4 \to 2 \to 0\))
        \item Minimalne rozszerzenie dla cyklu Hamiltona: \(0\)
        \item Maksymalna długość cyklu: \(5\)
    \end{itemize}
    \item Metryka podobieństwa pomiędzy grafem 1 i grafem 2: \(1.000\)
    \item Metryka podobieństwa pomiędzy grafem 2 i grafem 3: \(2.200\)
\end{itemize}




